\documentclass[12pt]{article}
\usepackage[a4paper,top=3cm,right=2cm,left=2cm,bottom=2.5cm]{geometry}
\usepackage[latin1]{inputenc}
\usepackage[T1]{fontenc}
\usepackage{newtxtext}
\usepackage[ttscale=0.875]{libertine}
\input{math.tex}
\usepackage[libertine,cmintegrals,cmbraces]{newtxmath}
\usepackage[round]{natbib}
\usepackage[american]{babel}
\usepackage[small,compact]{titlesec}
\usepackage{amsmath}
\usepackage{listings}
\usepackage{array}
\usepackage{tabularray}
\usepackage{diagbox}
\usepackage{comment}
\usepackage[defaultlines=4,all]{nowidow}
\usepackage{needspace}
\usepackage{graphicx}
\usepackage{mathrsfs}
\usepackage{ifpdf}
%\usepackage{upgreek}
\usepackage[pdftex,colorlinks=true,%
   urlcolor=blue,
   linkcolor=blue,
   citecolor=blue,
   filecolor=blue]{hyperref}
\usepackage{iexec}
\input{mylistings.tex}
\input{mydescription.tex}
\input{ira.tex}
\renewcommand{\ttdefault}{lmtt}% default cmtt
\newcommand{\texttk}[1]{\texttt{\bfseries #1}}
\newcommand{\ul}{\underline}

\usepackage{titlesec}

\UseTblrLibrary{booktabs}

% Customize the part title: [shape]{format}{label}{sep}{before}
\titleformat{\part}[display]
  {\normalfont\LARGE\bfseries\raggedright} % Font size (e.g., \huge, \Large, \large)
  {\partname\ \thepart}                 % Label (e.g., Part I)
  {20pt}                                % Separation between label and title
  {\LARGE}                               % Code before the title


\newcolumntype{C}[1]{>{\centering\arraybackslash}p{#1}}

\title{Documentation for \texttt{geo.chpl}}

\author{Nelson Lu�s Dias (\texttt{nelsonluisdias@gmail.com})}

\date{\today}

\begin{document}

\maketitle

\begin{abstract}
\texttt{geo} is a module for things that vary spatially
\end{abstract}


\part{Theory}\label{part:theory}

A vector $(u,v)$ can be in either of 4 quadrants. In this module, the
quadrants are 0,1,2,3 (instead of 1,2,3,4). In procedure \texttt{inxy}
below, we compute consecutive vectors from a fixed point $(x_p,y_p)$
to the points $(x_b,y_b)$ along a closed boundary. For each
consecutive pair, there is an increment that determines the difference
in quadrants between the first and second vectors.

These increments can be subsumed in a ``from/to'' matrix. The ``from''
is the quadrant of the first vector, and is listed along the lines $i$
of the matrix; the ``to'' is the quadrant of the second vector, and is
listed along the columns $j$. Most, but not all, of the matrix's
elements $(i,j)$ are equal
to $j-i$. In position $(0,3)$, $+3$ is
replaced by $-1$, because when the first vector is in quadrant 0 and the
second vector is in quadrant 3, there is a ``$-1$'' rotation; by the
same token, when the first vector is in quadrant 3 and the second vector
is in quadrant 0, there is a ``$+1$'' rotation. In quadrant units, a
rotation of $+1$ is equivalent to a rotation of $\uppi/2$ radians.

In \texttt{inxy}, we sum the increments as $(x_b,y_b)$ wanders around
the boundary. If $(x_p,y_p)$ is outside of the boundary, the sum will
be zero; otherwise, it will be $\pm 4$.

\newcommand{\scr}[1]{\scriptsize{#1}}

\begin{table}
  \caption{The skew-symmetric matrix of quadrant increments $j-i$. In position $(0,3)$, $+3$ is
    replaced by $-1$, and in position $(3,0)$, $-3$ is replaced by $+1$  \label{tab:from-to}}\centering
  \begin{tabular}{|l|C{1.2cm}|C{1.2cm}|C{1.2cm}|C{1.2cm}|}
    \hline
    \diagbox{from $i$}{to $j$} & 0              & 1              & 2              & 3              \\
    \hline
    0             & $\phantom{+}0$ & $+1$           & $+2$           & $-1$           \\
    1                          & $-1$           & $\phantom{+}0$ & $+1$           & $+2$           \\
    2                          & $-2$           & $-1$           & $\phantom{+}0$ & $+1$           \\
    3                          & $+1$           & $-2$           & $-1$           & $\phantom{+}0$ \\
    \hline
  \end{tabular}
\end{table}
  

\clearpage


\part{Constants and procedures}\label{part:const-procs}

\section{Constants}

There are no constants yet.



\section{\texttt{findquad}: in which quadrant is a point/vector?}\label{sec:findquad}

\subsection{Use}

\iexec[quiet,trace,stdout=xtract.txt]{grep 'private proc findquad' -A3
  geo.chpl}


\lstinputlisting[style=lstchapel, numbers=none,
  title={ \texttt{findquad}: finds the quadrant to which the point/vector
    (delx,dely) belongs.}
]{xtract.txt}

\subsection{Theory}

\texttt{findquad} is private and should not be called from outside of
\texttt{geo.chpl}. It is fairly self-explanatory, so check the source
code.

\section{\texttt{inxy}: is the point (xp,yp) inside perimeter (axboun,ayboun)?}\label{sec:inxy}

\subsection{Use}

\iexec[quiet,trace,stdout=xtract.txt]{grep 'proc inxy' -A5
  geo.chpl}


\lstinputlisting[style=lstchapel, numbers=none,
  title={ \texttt{inxy}: is the point (xp,yp) inside perimeter (axboun,ayboun)?}
]{xtract.txt}

\subsection{Theory}

\texttt{inxy} uses the ``Winding point algorithm'' (see
\url{https://en.wikipedia.org/wiki/Point_in_polygon}) and follows the
ideas in \cite{weiler1994incremental}, but my implementation is
considerably simpler than Weiler's (implemented in C).

Given the point $P = (x_p,y_p)$, \texttt{inxy} draws sucessive vectors
\[
(\delta_{xi},\delta_{yi}) = (x_{bi} - x_p,y_{bi} - y_p)
\]
from $P$ to each point ``$i$'' $(x_{bi},y_{bi})$ in the boundary, for $i=0,\ldots,n$. We
require that
\begin{align*}
  x_{bn} &= x_{b0}, \\
  y_{bn} &= y_{b0},
\end{align*}
so that the boundary ``closes''. We define the ``rotation'' from
$(\delta_{xi-1},\delta_{yi-1})$ to $(\delta_{xi},\delta_{yi})$ by the
quadrant increments of table \ref{tab:from-to}, and sum the rotations
\texttt{delquad} from 1 to $n$. If $P$ is an interior point, the total
sum of the rotations should be $\pm 4$; otherwise, it is 0.


\begin{comment}


\section{\texttt{viscair}: dynamic viscosity of air}

From \cite{montgomery--viscosity-air}.


\iexec[quiet,trace,stdout=xtract.txt]{grep 'proc viscair' -A2
  atmgas.chpl}

\lstinputlisting[style=lstchapel, numbers=none,
  title={ \texttt{viscair}: dynamic viscosity of air
    ($\mathrm{Pa\,s}$).}
]{xtract.txt}

\section{\texttt{difmom}: kinematic viscosity of air}

\iexec[quiet,trace,stdout=xtract.txt]{grep 'proc difmom' -A4
  atmgas.chpl}

\lstinputlisting[style=lstchapel, numbers=none,
  title={ \texttt{difmom}: kinematic viscosity of air
    ($\mathrm{m^2\,s^{-1}}$).}
]{xtract.txt}

\section{\texttt{difheat}: thermal diffusivity of air}

\iexec[quiet,trace,stdout=xtract.txt]{grep 'proc difheat' -A4
  atmgas.chpl}

\lstinputlisting[style=lstchapel, numbers=none,
  title={ \texttt{difheat}: thermal diffusivity of air
    ($\mathrm{m^2\,s^{-1}}$).}
]{xtract.txt}

\section{\texttt{difvap}: diffusivity of water vapor in air}

\iexec[quiet,trace,stdout=xtract.txt]{grep 'proc difvap' -A4
  atmgas.chpl}

\lstinputlisting[style=lstchapel, numbers=none,
  title={ \texttt{difvap}: thermal diffusivity of air
    ($\mathrm{m^2\,s^{-1}}$).}
]{xtract.txt}


\section{\texttt{rho2x}: convert density to molar fraction}

\iexec[quiet,trace,stdout=xtract.txt]{grep 'proc rho2x' -A5
  atmgas.chpl}

\lstinputlisting[style=lstchapel, numbers=none,
  title={ \texttt{rho2x}: convert density to molar fraction.}
]{xtract.txt}


\section{\texttt{pressalt}: atmospheric pressure as a function of altitude}

\iexec[quiet,trace,stdout=xtract.txt]{grep 'proc pressalt' -A2
  atmgas.chpl}

\lstinputlisting[style=lstchapel, numbers=none, title={
    \texttt{pressalt}: atmospheric pressure as a function of altitude
    in a standard atmosphere.}  ]{xtract.txt}

\end{comment}

\bibliography{abbrevs,all}
\bibliographystyle{apalike}
\end{document}


